\begin{topic}[id=starlink_megaconstellations,
             difficulty={Mittel},
             type={Systemanalyse + Bewertung},
             prerequisites={Grundlagen Kommunikationssysteme; Interesse an Netzinfrastrukturen},
             practice={optional},
             owner={Dozententeam},
             updated={2026-04-09},
             tags={ LEO, Mega-Konstellation, Starlink, Latenz, Kapazität, Bodenstationen }]{Starlink \& Mega-Konstellationen}

\TopicDescription{LEO-Konstellationen wie Starlink versprechen Breitband mit niedriger Latenz – auch außerhalb klassischer Netze. Wir beleuchten Architektur (Satelliten, Bodenstationen, Laserlinks), Kapazität/Abdeckung, Latenzpfade sowie Einschränkungen (Wetter, Sicht, Terminal). Ziel: nüchterne Einordnung der Leistungsfähigkeit und der Auswirkungen (Astronomie, Regulierung, Nachhaltigkeit).}

\TopicLeitfragen{\item Wie entsteht die niedrige Latenz im LEO?
\item Welche Kapazitätsgrenzen gibt es regional?
\item Welche Nebenwirkungen sind relevant (Astronomie, Weltraummüll)?
}
\TopicScope{\item Kein Tarif-/Providervergleich.
\item Keine Einzelhardware-Reviews.
}
\TopicPractice{Optionale Praxisidee: Pfadanalyse (Satellit/Backhaul) und Latenzabschätzung für ein Szenario.}

\TopicKeywords{LEO, Mega-Konstellation, Starlink, Latenz, Kapazität, Bodenstationen}

\nocite{spacexStarlink,nsfLaserLinks}
\end{topic}
